\documentclass[12pt, a4paper]{article}

% Paquetes requeridos
\usepackage[utf8]{inputenc}
\usepackage[T1]{fontenc}
\usepackage[spanish]{babel}
\usepackage{amsmath}
\usepackage{geometry}

% Configuracion de la pagina
\geometry{a4paper, margin=2.5cm}

\begin{document}

\section*{Análisis de corrientes ante un cortocircuito entre $a$ y $d$}

Si se presenta un cortocircuito directo entre los nodos $a$ y $d$, el potencial en ambos nodos se iguala ($V_a = V_d = V$), y como consecuencia del comportamiento de la fuente dependiente, el nodo $c$ también alcanza este mismo potencial ($V_c = V$). Al encontrarse los nodos $a$, $c$ y $d$ exactamente a la misma tensión, la diferencia de potencial a través de las resistencias $R_2$ y $R_3$ se vuelve nula. Por aplicación directa de la Ley de Ohm, la corriente que fluye por $R_2$ y $R_3$ será de $0\,\text{A}$.

\vspace{0.3cm}

Al no haber aporte ni sustracción de corriente por parte de las resistencias en el nodo $c$, la Ley de Corrientes de Kirchhoff (LKC) en dicho nodo fuerza a que la fuente dependiente de corriente tenga un valor nulo ($2i_x = 0$), lo que a su vez demuestra que $i_x = 0\,\text{A}$. De esta manera, la corriente a través de la fuente de tensión independiente $V$ deja de fluir. Toda la dinámica del circuito se simplifica y queda confinada: la corriente a través de la resistencia $R_1$ queda fijada de manera estática en $V/R_1$, y cualquier desbalance inyectado por la fuente de corriente $I$ fluirá exclusivamente a través del nodo $a$ y el cortocircuito hacia el nodo $d$, retornando por tierra para satisfacer las leyes de conservación de carga en el nodo de referencia $b$.

\end{document}